\begin{figure}[t]
\centering
\resizebox{\textwidth}{!}{
\begin{tikzpicture}[
box/.style={
rectangle,draw,rounded corners,
align=center,
minimum width=4.3cm,
minimum height=1.3cm
},
arrow/.style={->,thick}
]

% INPUTS
\node[box] (fit)
{Spectroscopic fit\\Watson constants\\$(D,H)$};

\node[box,below=1.6cm of fit] (harm)
{Harmonic input\\geometry + Hessian\\or `.fchk`};

\node[above=0.6cm of fit] {\large \textbf{Input}};

% LIFT
\node[box,right=4.9cm of fit] (lift)
{Lift / reconstruction\\pseudoinverse gauge fixing\\tensor representatives};

\node[above=0.6cm of lift] {\large \textbf{Lift}};

% TRANSPORT
\node[box,right=4.9cm of lift] (transport)
{Transport\\axis permutation\\reduction / representation change};

\node[above=0.6cm of transport] {\large \textbf{Transport}};

% OUTPUT
\node[box,right=4.9cm of transport,yshift=1.35cm] (reproj)
{Standard output\\reprojected Watson constants\\in chosen parametrization};

\node[box,right=4.9cm of transport,yshift=-1.35cm] (diag)
{Post-standard diagnostics\\conditioning, $s_{111}$, $H_{22}$,\\linear sextic $\tau(H_{22})$ term};

\node[above=0.6cm of reproj] {\large \textbf{Output}};

% ARROWS
\draw[arrow] (fit.east) -- (lift.west);
\draw[arrow] (harm.east) -- (diag.west);
\draw[arrow] (lift.east) -- (transport.west);
\draw[arrow] (transport.east) -- (reproj.west);
\draw[arrow] (transport.east) -- (diag.west);

\end{tikzpicture}
}

\caption{
Operational structure of the present framework. Standard quartic and sextic
centrifugal distortion constants are lifted from Watson parametrizations to
tensor representatives through pseudoinverse gauge fixing, transported under
representation changes, and reprojected onto the target Hamiltonian
parametrization. Optional harmonic input is only needed for post-standard
diagnostics, namely the quartic $H_{22}$ contribution and its first sextic
linear analogue.
}

\label{fig:ceditt3_scheme}
\end{figure}
